\section{Comparison With a Deep-Learning Baseline}\label{sec:cnn}

A common approach in the CASIA~v2.0 literature feeds error-level-analysis (ELA) images to a convolutional network and
reports accuracies around 90\%~\cite{sudiatmika2019ela,ali2022recompress}. We built such a baseline, with the same
splits, protocols and frozen test set, to see whether a learned model beats the hand-crafted pipeline once the format
leak is controlled, and whether those high numbers survive that control.

\subsection{Method}\label{sec:cnn-method}

ELA is computed on the image \emph{after} the protocol: the absolute per-channel difference to its JPEG re-save at
quality~90, scaled by 10, clipped to 8~bits and kept at native resolution (examples in Appendix~\ref{app:cnn}). A
ResNet-18~\cite{he2016resnet} with ImageNet weights~\cite{deng2009imagenet} and a single logit is trained on one random
$192\times192$ crop per image per epoch with class-weighted binary cross-entropy and AdamW (details in
Appendix~\ref{app:cnn}); images smaller than a tile are zero-padded, which in ELA space means ``no error''. The image
score pools the logits of overlapping $192\times192$ tiles (stride 96) by their mean or maximum. Only the epoch and the
pooling rule were chosen, by validation AUC; the recipe was then refit on all 10,722 development images, as was the
classical forest. Under R, three seeds differ little on validation (SD 0.008), so one run is representative. Before
the network saw the test split, the comparisons were pre-registered with the SHA-1 of every model; primary were the R
test AUC, the paired difference CNN $-$ classical under R and B, and a fixed rank-average combination of the two
scores. The evaluation ran once. The comparator is the forest on the forensic features (local + global) of
Section~\ref{sec:test}.

\begin{table*}[!t]
\centering
\caption{ELA-CNN vs.\ the classical pipeline: test ROC-AUC [95\% CI] and paired differences (group bootstrap,
1,000 resamples). Classical = random forest on the forensic features.}
\label{tab:cnn-test}
\footnotesize
\begin{tabular}{lccc}
\toprule
 & R & B & A \\
\midrule
ELA-CNN & 0.770\ci{0.744}{0.795} & 0.811\ci{0.788}{0.834} & \textbf{0.994}\ci{0.991}{0.997} \\
Classical & \textbf{0.805}\ci{0.779}{0.831} & \textbf{0.920}\ci{0.905}{0.934} & 0.981\ci{0.975}{0.986} \\
Rank average (CNN, classical) & \textbf{0.824}\ci{0.803}{0.846} & 0.903\ci{0.886}{0.919} & 0.996\ci{0.994}{0.998} \\
\midrule
CNN $-$ classical & $-$0.036\ci{$-$0.067}{$-$0.005} & $-$0.109\ci{$-$0.128}{$-$0.088} & +0.014\ci{0.008}{0.019} \\
CNN $-$ shortcut features (RF) & +0.058\ci{0.030}{0.084} & $-$0.075\ci{$-$0.094}{$-$0.055} & $-$0.005\ci{$-$0.008}{$-$0.003} \\
CNN $-$ shortcut features (SVM) & +0.017\ci{$-$0.011}{0.044} & $-$0.072\ci{$-$0.091}{$-$0.053} & -- \\
Rank average $-$ classical & +0.019\ci{0.003}{0.035} & $-$0.017\ci{$-$0.025}{$-$0.008} & +0.015\ci{0.011}{0.019} \\
\bottomrule
\end{tabular}
\end{table*}

\begin{figure*}[!t]
\centering
\includegraphics[width=\textwidth]{fig_cnn_roc.png}
\caption{Test ROC curves of the ELA-CNN, the classical forest, their rank average and the shortcut features under
protocols R, B and A (AUC in the legends). Only on the files as released (A) does the CNN reach a very high AUC (0.994),
and there the shortcut features alone reach 1.000.}
\label{fig:cnn-roc}
\end{figure*}

\subsection{Test results}\label{sec:cnn-test}

Table~\ref{tab:cnn-test} and Fig.~\ref{fig:cnn-roc} give the results. With a Holm correction~\cite{holm1979} over the
four primary paired comparisons (CNN $-$ classical and rank average $-$ classical, under R and B), both R results stay
significant (adjusted $p=0.036$ each; unadjusted 0.026 and 0.018, bootstrap $p$-values); the B results have $p<0.002$.
The two R differences are significant on this test set but small, and not adjusted for variation between training
runs, which the test CIs do not include.

\emph{Under the strict protocols the classical pipeline wins.} Under R the CNN is 0.036 below the forensic forest. It
beats the shortcut features' random forest (+0.058) but is not significantly better than their SVM (0.753). Since the
CNN can also see global cues such as the overall ELA level and, through max pooling, image size, we cannot show that
it learns more than the global shortcuts do~\cite{geirhos2020shortcut}. Under B the gap widens to 0.109 and the CNN
falls below B's global shortcuts (0.811 vs.\ 0.886): B keeps within-image compression traces that the DCT, ghost
and ELA modules measure explicitly, and the CNN exploits them less well.

\emph{On the files as released (A)} the CNN reaches AUC 0.994 (the literature reports ${\sim}$90\% accuracy on these
files), while the shortcut features alone reach 1.000 on the same files: the leak alone gives such
a score. It is visible in ELA itself (Fig.~\ref{fig:cnn-ela}, Appendix~\ref{app:cnn}) and persists on JPEG files only
(0.995). A plausible explanation, not tested directly: authentic JPEGs mostly use the standard IJG
tables (89\%), tampered ones never do (Section~\ref{sec:dataset}), and ELA is sensitive to the table. This is
incomplete, since the CNN also separates tampered JPEGs from the 126 authentic JPEGs with non-standard tables (AUC
0.971).

\emph{Under R the two approaches are complementary.} Their scores are moderately correlated (Spearman 0.57), and their
rank average reaches 0.824, the highest R AUC among detectors without shortcut features, level with the best model of
Section~\ref{sec:test}, the SVM on forensic + shortcut features (0.829; difference $-$0.005\ci{$-$0.024}{0.014}). Under
B, averaging in the weaker CNN hurts. The combination ranks scores within the evaluation set and cannot score a single
image, so the deployed detector remains the classical one.

\begin{table}[!t]
\centering
\caption{Test AUC under R by forgery type (authentic vs.\ one type) and by natural image size (post-hoc).}
\label{tab:cnn-groups}
\scriptsize
\setlength{\tabcolsep}{4pt}
\begin{tabular}{lccc}
\toprule
 & CNN & Classical & CNN $-$ classical \\
\midrule
Copy-move & 0.778 & 0.874 & $-$0.096\ci{$-$0.125}{$-$0.069} \\
Splicing & 0.755 & 0.682 & +0.073\ci{0.019}{0.128} \\
\midrule
Smaller than $384\times256$ ($n$=21) & 0.702 & 0.615 & +0.087\ci{$-$0.072}{0.260} \\
$384\times256$ ($n$=1,469) & 0.770 & 0.789 & $-$0.019\ci{$-$0.065}{0.026} \\
Larger than $384\times256$ ($n$=402) & \textbf{0.544} & 0.730 & $-$0.186\ci{$-$0.232}{$-$0.139} \\
\bottomrule
\end{tabular}
\end{table}

\subsection{Where each approach wins}\label{sec:cnn-groups}

Table~\ref{tab:cnn-groups} splits the R results. Keypoint matching is much better at copy-move than the CNN, which is
not built to compare distant regions. On splicing, the classical pipeline's main weakness (Section~\ref{sec:test}), the
CNN wins (+0.073), but the global shortcut features reach almost the same AUC (SVM 0.743; forensic + shortcut SVM
0.752), so this edge may come from global compression or size differences between CASIA's spliced and authentic
images rather than from local splice traces. It is still what makes the rank average help.

The pre-registered size check (AUCs within area terciles) is \emph{invalid}: 78\% of the test images are
$384\times256$ in either orientation, and the code broke this tie by row order, which lists authentic images first, so
the terciles were confounded with the label. A post-hoc grouping by natural size replaces it. On standard-size images
the CNN is not significantly different from the forest, though the wide CI does not show equivalence; most of its
deficit comes from larger images, where it is near chance (0.544); the forest, some of whose features also depend on
size, drops less. Larger authentic images receive much higher CNN scores than standard-size ones (mean logit 1.35
vs.\ $-$0.40): the network appears to have learned the ELA statistics of the dominant $384\times256$ images and to
treat others as suspicious (an interpretation; the image sources were not analysed).

\emph{Transfer.} The CNN does not transfer to MICC-F220: AUC 0.529\ci{0.360}{0.700} under R, consistent with chance,
against 0.972 for the classical detector (11 forged scenes, hence the wide CI). Degradation results are in
Appendix~\ref{app:cnn}.

\subsection{Summary and limitations}\label{sec:cnn-summary}

A standard ELA-CNN reaches high scores only on the leaking files. Under R and B it is below the
classical pipeline and not significantly better than a shortcut-feature SVM, but under R its rank average with the
forest matches the best model. Its failure on MICC-F220 is the strongest argument in this study for explicitly modelled
classical cues. Limitations: one architecture and input representation, no search beyond epoch and pooling, a single
validation split, and non-bit-reproducible GPU training. A stronger network, e.g.\ trained end-to-end on RGB with
constrained convolutions~\cite{bayar2016constrained}, might close the gap under R, but would need the same leak
controls.
